% ---------------------------------------------------------
% Theorems
% ---------------------------------------------------------

\newtheorem{theorem}{Theorem}[chapter]
\newtheorem{lemma}[theorem]{Lemma}
\newtheorem{proposition}[theorem]{Proposition}
\newtheorem*{corollary}{Corollary}
\newtheorem{definition}{Definition}[chapter]
\newtheorem{conjecture}{Conjecture}[chapter]
\newtheorem{example}{Example}[chapter]
\newtheorem{assumption}{Assumption}[chapter]

\theoremstyle{remark}
\newtheorem*{remark}{Remark}
\newtheorem*{note}{Note}
\newtheorem{case}{Case}

% ---------------------------------------------------------
% Graphic paths
% ---------------------------------------------------------

\DeclareGraphicsExtensions{.pdf,.png,.jpg}
\newcommand{\figurepath}{{images/}}
\graphicspath{{images/}}

% ---------------------------------------------------------
% Spacing
% ---------------------------------------------------------

\renewcommand{\arraystretch}{1.15}

\setlength{\parindent}{0pt}
\setlength{\parskip}{1.0ex plus 0.0ex minus 0.0ex}

\setcounter{tocdepth}{1}


% ---------------------------------------------------------
% Counter
% ---------------------------------------------------------
\counterwithout{footnote}{chapter} 


% ---------------------------------------------------------
% Listings
% ---------------------------------------------------------

\renewcommand{\lstlistlistingname}{Verzeichnis der Listings}
\definecolor{darkgreen}{rgb}{0,.4,0}
\definecolor{darkviolett}{rgb}{.4,0,.4}
\definecolor{blue}{rgb}{0,0,.9}
\newcommand{\listingswidth}{\textwidth-1cm}
\lstset{
    language=Java,                  % oder  C++, Pascal, {[77]Fortran}, ...
    numbers=left,                   % Position der Zeilennummerierung
    firstnumber=auto,               % Erste  Zeilennummer
    basicstyle=\ttfamily\footnotesize,     % Textgröße  des Standardtexts
    keywordstyle=\ttfamily\bfseries\color{blue},    % Formattierung Schlüsselwörter
    commentstyle=\ttfamily\color{gray},       % Formattierung Kommentar
    stringstyle=\ttfamily\color{darkgreen},             % Formattierung Strings
    numberstyle=\small,              % Textgröße der Zeilennummern
    stepnumber=1,                   % Angezeigte Zeilennummern
    numbersep=5pt,                  % Abstand zw. Zeilennummern und Code
    aboveskip=15pt,                 % Abstand oberhalb des Codes
    belowskip=11pt,                 % Abstand unterhalb des Codes
    captionpos=b,                   % Position der Überschrift
    linewidth=\textwidth,               % \textwidth-2em
    xleftmargin=10pt,               % Linke Einrückung
    frame=lines,                   % Rahmentyp
    breaklines=true,                % Umbruch langer Zeilen
    showstringspaces=true,          % Spezielles Zeichen für Leerzeichen
    tabsize=2,
    escapeinside={@}{@)}			% Escaping labels 
}

% ---------------------------------------------------------
% Links and Hyperlinks
% ---------------------------------------------------------

\hypersetup{
    breaklinks=true,           % Zeilenumbruch bei Links
    colorlinks=true,           % true = farbige Links, false = Rahmen
    linkcolor=blue,            % Farbe für Links auf gleicher Seite
    citecolor=blue,            % Farbe für Links auf Zitatstellen
    filecolor=blue,
    urlcolor=blue,             % Farbe für Links (nicht nur URLs)
    anchorcolor=blue,          % Farbe für Anker (?)
    menucolor=red,             % Farbe für Acrobatmenüeinträge
    plainpages=false           % Fehlermeldungen reduzieren
}
